\section{Model Agnostic: sin religión de modelos, pero con disciplina de evaluación}

A partir de los casos de la industria, el ponente plantea la postura model agnostic: no construir una identidad en torno a open/closed, small/large ni a un vendor en particular, sino elegir el modelo según los objetivos de salvar vidas, eficiencia y productividad. Esta postura no sostiene que «todos los modelos son iguales»; exige establecer task contract, evaluation, routing, fallback y version governance más rigurosos.

\subsection{Portfolio selection en lugar de un winner único}

Esta sección convierte «no model religion» en un proceso de selección continuo, y no en un lema de neutralidad ante los proveedores. Al leer la figura, se parte del task contract y después se comparan los candidatos, la evidencia y la routing policy; así, aunque las capacidades y los precios de los modelos cambien con rapidez, el sistema puede explicar por qué cierto tipo de solicitud usa cierto modelo, cuándo se recurre al fallback, cómo se revierte el cambio y qué restricciones de gobernanza no pueden quedar anuladas por una mejora de rendimiento.

\lecturefigure{10-model-agnostic.png}{La estrategia model-agnostic tiene como eje la tarea, la combinación de candidatos, la evidencia y la routing policy.}{Subtítulos locales 20:20--23:40; figura conceptual redibujada.}

\begin{knowledgebox}{Lectura de la figura: cuatro pasos para evitar la model religion}
Primero se definen la calidad de la tarea, la latency, la privacy, el language y el cost; después se forman candidatos open/closed, small/large y general/specialist; luego se ejecutan offline/online eval y red-team; por último, la elección se escribe como una routing policy con capacidad de auditoría. Cuando llega un modelo nuevo, se vuelve a comparar, en lugar de sustituir automáticamente todas las rutas.
\end{knowledgebox}

\begin{importantbox}{La selección de modelos es una optimización con restricciones}
Puede escribirse como
\[
m^*=\arg\max_m Q(m)\quad
\text{s.t.}\quad L(m)\leq L_{max},\ C(m)\leq C_{max},\ G(m)=1,
\]
donde $Q$ es la calidad de la tarea, $L$ es la latency, $C$ es el costo y $G$ indica si se cumplen las restricciones de gobernanza, de datos y de despliegue. Si la distribución de tareas varía, la solución óptima puede ser routing o ensemble, y no un único modelo.
\end{importantbox}

\begin{warningbox}{Model agnostic no equivale a independencia de la cadena de suministro}
El modelo puede reemplazarse, pero el tokenizer, el tool schema, la safety policy, la observabilidad, los datos de fine-tuning y el evaluation history no necesariamente son portables. La verdadera portability requiere versioned interface, prompt contract, data export, fallback y rollback, no agregar un menú desplegable a la UI.
\end{warningbox}

\teachervoice{El ponente extiende «no technology religion» a «no model religion». La versión de ingeniería de esta frase es: mantener la arquitectura falsable, mantener al proveedor reemplazable y conservar evidencia comparable de cada migración.}

\subsection{Resumen de la sección}

Model agnostic es un portfolio guiado por evaluation, no la idea de que «cualquier modelo sirve». El contrato de la tarea, las restricciones, las versiones y el routing determinan juntos la elección; cuanto más a menudo se cambia de modelo, más se necesitan interfaces de sistema estables y pruebas de regresión.
